\section{Метод решения}
Решение разделено на общую логику лабораторной работы №2, платформенный слой процессов и новый модуль \texttt{ipc}. Родительская программа читает ввод пользователя, дочерняя проверяет числа и записывает результат. Модуль \texttt{ipc} создаёт отображаемый файл и объекты уведомления; он скрывает различия Windows и Linux за классом \texttt{ipc::Channel}.

\subsection{Отображаемый файл}
Родитель создаёт временный файл размером с общую структуру и отображает его в своё адресное пространство. Дочерний процесс получает путь к этому файлу аргументом командной строки, открывает его и отображает у себя. Обе программы видят одни и те же страницы памяти, поэтому запись одного процесса становится видна другому без передачи данных через ядро~\cite{linuxmmap,winsharing,winmapping}.

\begin{center}
\begin{tabular}{lll}
Действие & Linux & Windows \\
\hline
Создание файла & \texttt{mkstemp}, \texttt{ftruncate} & \texttt{GetTempFileNameW}, \texttt{CreateFileW} \\
Отображение & \texttt{mmap(MAP\_SHARED)} & \texttt{CreateFileMappingW}, \texttt{MapViewOfFile} \\
Открытие дочерним & \texttt{open}, \texttt{mmap} & \texttt{CreateFileW}, \texttt{MapViewOfFile} \\
Удаление файла & \texttt{unlink} в родителе & \texttt{FILE\_FLAG\_DELETE\_ON\_CLOSE} \\
\end{tabular}
\end{center}

В Windows временный файл открывается с флагом удаления при закрытии и режимом разделения \texttt{FILE\_SHARE\_DELETE}. Флаг действует после закрытия последнего дескриптора файла, поэтому дочерний процесс может открыть тот же файл, пока родитель его не закрыл~\cite{deleteonclose}.

\subsection{Синхронизация}
Общая память не уведомляет другой процесс о появлении данных. Поэтому каждое направление обмена снабжено собственным объектом ожидания: <<запрос готов>> для передачи от родителя к ребёнку и <<ответ готов>> для обратного направления. Родитель всегда записывает команду, подаёт сигнал запроса и ждёт ответа. Ребёнок ждёт запроса, обрабатывает его, записывает ответ и подаёт сигнал. В каждый момент записывает в общую структуру только один процесс, поэтому дополнительные блокировки не нужны.

В Linux используются неименованные семафоры POSIX. Они размещены в самом отображаемом файле и инициализируются вызовом \texttt{sem\_init} с ненулевым параметром \texttt{pshared}, что позволяет использовать их из разных процессов~\cite{linuxsem,linuxsemoverview}. В Windows используются именованные события (\texttt{CreateEventW} в родителе и \texttt{OpenEventW} в ребёнке) с автоматическим сбросом. Имя события строится из уникального имени временного файла~\cite{winevent}.

\subsection{Протокол обмена}
Общая структура \texttt{ipc::Block} содержит команду, ответ, идентификатор родителя, длину строки, признак переполнения, буфер строки на 256 байт и буфер текста ошибки на 256 байт.

\begin{center}
\begin{tabular}{ll}
Команда (родитель) & Значение \\
\hline
\texttt{Line} & В буфере лежит строка с числом \\
\texttt{Stop} & Конец ввода: ребёнку следует завершиться \\
\end{tabular}
\end{center}

\begin{center}
\begin{tabular}{ll}
Ответ (ребёнок) & Значение \\
\hline
\texttt{Ready} & Выходной файл открыт, ребёнок готов \\
\texttt{Written} & Составное число записано в файл \\
\texttt{Ignored} & Число 0 или 1 пропущено \\
\texttt{Invalid} & Ошибка формата, выход за диапазон \texttt{int} или слишком длинная строка \\
\texttt{Stop} & Получено простое или отрицательное число \\
\texttt{Stopped} & Подтверждение команды \texttt{Stop} \\
\texttt{Error} & Системная ошибка в ребёнке, текст записан в \texttt{errorText} \\
\end{tabular}
\end{center}

Ребёнок сначала открывает выходной файл и только после этого подаёт ответ \texttt{Ready}. Родитель не принимает числа до его получения, поэтому ошибка открытия файла обнаруживается до начала обработки. После этого на каждую строку приходит ровно один ответ.

\subsection{Завершение и аварийные ситуации}
Ожидание сигнала выполняется с периодом 100~мс. При каждом истечении периода проверяется, жив ли второй процесс. Родитель вызывает \texttt{os::IsProcRunning}. Ребёнок проверяет, существует ли родитель с идентификатором, записанным в общую структуру: в Linux сравнивается результат \texttt{getppid}, в Windows открывается дескриптор процесса вызовом \texttt{OpenProcess}. Без такой проверки гибель одного процесса оставила бы второй в бесконечном ожидании. Перед признанием партнёра погибшим выполняется ещё одна неблокирующая попытка получить сигнал, поскольку ответ мог прийти непосредственно перед завершением.

Общая память освобождается родителем после завершения ребёнка. Это исключает уничтожение семафора, пока второй процесс ещё может обратиться к нему.

\section{Описание программы}
\subsection{Структура файлов}
Пути приведены относительно корня проекта \texttt{Operation-Systems}.

\begin{itemize}
\item \texttt{lab4/code/parent/main.cpp} --- ввод имени файла и команд, запуск ребёнка, обмен через канал, ожидание завершения.
\item \texttt{lab4/code/child/main.cpp} --- обработка команд, проверка чисел и запись результатов.
\item \texttt{lab4/code/include/ipc.h} --- общая структура, перечисления команд и ответов, интерфейс канала.
\item \texttt{lab4/code/ipc\_posix.cpp} --- реализация канала для Linux.
\item \texttt{lab4/code/ipc\_win.cpp} --- реализация канала для Windows.
\item \texttt{lab2/code/common.cpp} --- разбор числа и проверка простоты (общая логика).
\item \texttt{common/code/os/} --- запуск процессов и запись файла; добавлена функция \texttt{os::IsProcRunning}.
\end{itemize}

\subsection{Основные типы и функции}
\texttt{ipc::Channel::Create} вызывается родителем, а \texttt{ipc::Channel::Open} --- ребёнком. Метод \texttt{block} возвращает ссылку на общую структуру. \texttt{SignalRequest} и \texttt{SignalResponse} подают сигналы, а \texttt{WaitRequest} и \texttt{WaitResponse} ожидают их и возвращают \texttt{false}, если партнёр погиб. Деструктор канала освобождает отображение, объекты синхронизации и временный файл.

\texttt{os::IsProcRunning} --- неблокирующая проверка состояния процесса. В Linux она вызывает \texttt{waitpid} с флагом \texttt{WNOHANG}, и полученный код завершения сохраняется для последующего \texttt{os::WaitProc}. В Windows вызывается \texttt{WaitForSingleObject} с нулевым таймаутом. Остальные функции платформенного слоя не изменялись.

\subsection{Используемые средства Linux}
\begin{itemize}
\item \texttt{mkstemp}, \texttt{ftruncate}, \texttt{open}, \texttt{fstat}, \texttt{unlink} --- создание, проверка размера и удаление файла.
\item \texttt{mmap}, \texttt{munmap} --- отображение файла в память.
\item \texttt{sem\_init}, \texttt{sem\_post}, \texttt{sem\_timedwait}, \texttt{sem\_trywait}, \texttt{sem\_destroy} --- уведомления между процессами.
\item \texttt{fork}, \texttt{execv}, \texttt{waitpid}, \texttt{getppid} --- запуск, ожидание и проверка живости процессов.
\end{itemize}

\subsection{Используемые средства Windows}
\begin{itemize}
\item \texttt{GetTempPathW}, \texttt{GetTempFileNameW}, \texttt{CreateFileW} --- создание временного файла.
\item \texttt{CreateFileMappingW}, \texttt{MapViewOfFile}, \texttt{UnmapViewOfFile} --- отображение файла в память.
\item \texttt{CreateEventW}, \texttt{OpenEventW}, \texttt{SetEvent}, \texttt{WaitForSingleObject} --- события и ожидание.
\item \texttt{OpenProcess}, \texttt{CreateProcessW} --- проверка родителя и запуск ребёнка.
\item \texttt{MultiByteToWideChar}, \texttt{WideCharToMultiByte} --- преобразование путей. Путь к временной папке может содержать русские буквы, поэтому ребёнок использует точку входа \texttt{wmain}.
\end{itemize}

\subsection{Обработка ошибок}
Ошибки системных вызовов преобразуются в исключения с названием операции и кодом ОС. Внешние обработчики выводят диагностику в \texttt{stderr} и возвращают ненулевой код. Если ребёнок не смог записать число в файл, он помещает текст ошибки в общую структуру, отвечает \texttt{Error} и завершается с кодом 1. Если ребёнок не смог открыть выходной файл, он завершается до ответа \texttt{Ready}, а родитель обнаруживает это при ожидании и также возвращает ошибку. При любом исключении в родителе работающий ребёнок принудительно завершается.

\section{Исходный код}
Листинги подключаются непосредственно из рабочих исходников; отдельные копии кода в отчёте не хранятся.

\subsection{Общий интерфейс канала}
\lstinputlisting[language=C++,caption={Структура общей памяти и интерфейс канала},captionpos=b]{\codedir/include/ipc.h}

\subsection{Реализация для Linux}
\lstinputlisting[language=C++,caption={Отображаемый файл и семафоры POSIX},captionpos=b]{\codedir/ipc_posix.cpp}

\subsection{Реализация для Windows}
\lstinputlisting[language=C++,caption={Отображаемый файл и именованные события},captionpos=b]{\codedir/ipc_win.cpp}

\subsection{Родительский процесс}
\lstinputlisting[language=C++,caption={Ввод, запуск ребёнка и обмен через канал},captionpos=b]{\codedir/parent/main.cpp}

\subsection{Дочерний процесс}
\lstinputlisting[language=C++,caption={Обработка команд и запись составных чисел},captionpos=b]{\codedir/child/main.cpp}
